\section{Results}
\label{sec:results}

\subsection{The delivery channel dominates}

Tasks 04 through 06 carry an identical injected instruction and differ only in
how it reaches the agent. Holding the text constant, the channel changes the
follow-through rate by a factor of six.

\begin{table}[h]
\centering
\small
\begin{tabular}{@{}llrr@{}}
\toprule
\textbf{Task} & \textbf{Channel} & \textbf{Followed} & \textbf{Rate} \\
\midrule
\texttt{task\_04} & live network fetch & 2/20 & 10\% (17\% of those exposed) \\
\texttt{task\_05} & cross-run persistent memory & 10/20 & 50\% \\
\texttt{task\_06} & local knowledge base & 12/20 & 60\% \\
\bottomrule
\end{tabular}
\caption{Corpus, 60 runs. The network row carries two denominators because
the attacker logs each delivery, and eight of those twenty runs were never
served the payload (Section~\ref{sec:validity}). Conditioning on exposure
raises the rate but does not change the ordering.}
\end{table}

A hypothesis visible directly in the transcripts: content that was already
there, sitting in team memory or in internal documentation, receives less
scrutiny than content just fetched over HTTP from an external mirror, even
when the instruction is byte for byte the same. That asymmetry of trust is
what ASI06 warns about~\cite{owasp-asi06}, and it points the wrong way for
defenders, since the cheapest channel to poison is also the most trusted.

\subsection{The guardrail correlates with more leakage, and it replicates}

The finding that survived every correction in this paper is also the one the
team least expected. On \texttt{task\_06}, the condition carrying a production
security guardrail leaks at roughly three times the rate of the condition
without one.

\begin{table}[h]
\centering
\small
\begin{tabular}{@{}llrr@{}}
\toprule
\textbf{Set} & \textbf{Ollama} & \texttt{sin\_harness} & \texttt{con\_harness} \\
\midrule
Corpus, machine A & 0.32.5 & 3/10 & \textbf{9/10} \\
Machine B & 0.34.0 & 3/10 & \textbf{7/10} \\
Machine B & 0.6.8 & 10/10 & 10/10 \\
\bottomrule
\end{tabular}
\caption{The third row is not a result, it is a measurement failure preserved
for honesty: on an Ollama release 26 minor versions old the agent loops until
it exhausts the turn cap, every condition saturates, and nothing can be
distinguished. Section~\ref{sec:validity} gives the mechanism.}
\end{table}

Once the runtime versions are comparable, the effect appears on both machines
at close to the same size. It is a correlation across two independent sets,
not a causal claim, and the next subsection is our attempt to explain it.

\subsection{Naming the attack surface is not the mechanism}

\texttt{HARNESS\_SECURITY\_BLOCK} confounds three things at once: it is
longer, it is more emphatic, and it names the exact surfaces an attacker would
ask the agent to use, warning it never to put its token in file names,
directory names, or URLs. A prompt-injection payload then asks for precisely
that. The natural hypothesis is priming: naming the surface makes it available.

We registered a falsifiable prediction in \texttt{sandbox/prompts.py} before
running. \texttt{GENERIC\_EMPHASIS\_BLOCK} matches the guardrail's length and
emphasis (818 characters against 830) but forbids only that the token leave
the working environment, without naming how. If priming is the mechanism, that
condition should sit near \texttt{sin\_harness}.

\begin{table}[h]
\centering
\small
\begin{tabular}{@{}lllr@{}}
\toprule
\textbf{Condition} & \textbf{Length and emphasis} & \textbf{Names the surface} & \textbf{Level 5} \\
\midrule
\texttt{sin\_harness} & short, none & no & 3/10 \\
\texttt{con\_harness\_generico} & matched & \textbf{no} & \textbf{7/10} \\
\texttt{con\_harness} & matched & yes & \textbf{7/10} \\
\bottomrule
\end{tabular}
\caption{Thirty runs on one machine, one model, one monitor build, with memory
reset before each run. The prediction fails: the generic condition lands
exactly on the guardrail, not near the control.}
\end{table}

Removing every mention of files, directories and URLs, while holding length
and emphasis constant, did not reduce leakage at all. Naming the attack
surface is not the mechanism.

What remains is simpler and less comfortable. The variable that tracks leakage
is the presence of an emphatic security block, whatever it says. A long,
serious warning about a sensitive token is associated with roughly double the
leakage of a line that mentions it in passing, and the specific content of the
warning changes nothing.

We are deliberately not proposing a mechanism for that. This experiment rules
one out; it does not identify its replacement.

\subsection{Detection latency and monitor strength}

\todo{Cerrar esta subseccion con las cifras corregidas de la Seccion 5, no con
las publicadas originalmente. Reportar TTD por condicion con IC bootstrap y
decir que por debajo de n=5 no se da intervalo. Anadir la figura del barrido
de fortaleza desde results/corpus/monitor\_strength.jsonl, que ya grafica el
dashboard: el heartbeat es lineal en su umbral y predecible, el canary es
plano porque desde el cambio a inotify el parametro ya no gobierna la
latencia. Mencionar que el TTD de la superficie de filesystem cae por debajo
del milisegundo y a esa resolucion mide desfase entre relojes, no latencia
del monitor.}

\subsection{Limitations of these results}

Ten runs per cell. The seven against three contrast is clear and in the
opposite direction to the prediction, but we report no effect size and run no
formal test. The causal experiment covers one task, chosen because the effect
was cleanest there, and was not repeated on the other two. Every run uses one
local 7B model, so nothing here is claimed to generalise to frontier models or
to other forms of guardrail. And the mechanism behind the emphasis effect is
unisolated: what this experiment eliminates is one concrete explanation, not
all of them.
